\begin{apartats}

\begin{tria}{parametres-tasca-a}
\itemtria{original}{1,5}{2,5}
Determina els valors de $a$ i $b$ perquè la funció següent sigui contínua a tot $\mathbb{R}$.
\[
f(x)=\begin{cases}
  e^{x}+a & \si{x\le 0},\\[3pt]
  ax+b & \si{0<x<2},\\[3pt]
  4+\ln(x-1) & \si{x\ge 2}.
\end{cases}
\]

\begin{solucio}
Cada branca és contínua al seu interval: $e^x+a$ i $ax+b$ ho són a tot $\mathbb{R}$, i
$4+\ln(x-1)$ ho és per a $x>1$, en particular a $[2,+\infty)$. Només cal estudiar els
enganxaments.\\
En $x=0$: $f(0)=e^0+a=1+a$ i $\lim_{x\to0^+}f(x)=b$. Cal $b=1+a$.\\
En $x=2$: $\lim_{x\to2^-}f(x)=2a+b$ i $f(2)=4+\ln1=4$. Cal $2a+b=4$.\\
Substituint: $2a+1+a=4$, d'on $\boxed{a=1}$ i $\boxed{b=2}$.
\end{solucio}

\itemtria{un-parametre}{1,5}{2,5}
Determina els valors de $m$ perquè la funció següent sigui contínua a tot $\mathbb{R}$.
\[
g(x)=\begin{cases}
  3e^{x}+1 & \si{x\le 0},\\[3pt]
  (x-m)^2 & \si{0<x<1},\\[6pt]
  \dfrac{9}{x} & \si{x\ge 1}.
\end{cases}
\]

\begin{solucio}
Cada branca és contínua al seu interval: $3e^x+1$ i $(x-m)^2$ ho són a tot $\mathbb{R}$, i
$\dfrac9x$ ho és per a $x\ne0$, en particular a $[1,+\infty)$. Només cal estudiar els
enganxaments.\\
En $x=0$: $g(0)=3e^0+1=4$ i $\lim_{x\to0^+}g(x)=m^2$. Cal $m^2=4$: $m=2$ o $m=-2$.\\
En $x=1$: $\lim_{x\to1^-}g(x)=(1-m)^2$ i $g(1)=9$. Cal $(1-m)^2=9$: $m=-2$ o $m=4$.\\
Només $\boxed{m=-2}$ compleix totes dues condicions. Amb $m=2$ seria contínua en $x=0$, però
en $x=1$ tindria un salt finit ($1$ per l'esquerra i $9$ per la dreta).
\end{solucio}
\end{tria}

\begin{nomesllarg}
\apartat{1}
Troba \textbf{tots} els valors de $m$ per als quals la funció següent és contínua a tot
$\mathbb{R}$.
\[
h(x)=\begin{cases}
  x+m^2 & \si{x<1},\\[3pt]
  x^2+3m & \si{x\ge 1}.
\end{cases}
\]

\begin{solucio}
Les dues branques són polinòmiques, i per tant contínues: només cal estudiar $x=1$.\\
$\lim_{x\to1^-}h(x)=1+m^2$ i $h(1)=1+3m$.\\
Cal $1+m^2=1+3m$, és a dir $m^2-3m=m(m-3)=0$: $\boxed{m=0}$ o $\boxed{m=3}$.
Hi ha \textbf{dos} valors. (Comprovació amb $m=3$: $1+9=10$ i $1+9=10$.)
\end{solucio}
\end{nomesllarg}

\end{apartats}
